
\documentclass{article}
\usepackage{colortbl}
\usepackage{makecell}
\usepackage{multirow}
\usepackage{supertabular}

\begin{document}

\newcounter{utterance}

\centering \large Interaction Transcript for game `cladder', experiment `full\_v1.5\_default', episode 7517 with selfrefine{-}Qwen3.5{-}27B.
\vspace{24pt}

{ \footnotesize  \setcounter{utterance}{1}
\setlength{\tabcolsep}{0pt}
\begin{supertabular}{c@{$\;$}|p{.15\linewidth}@{}p{.15\linewidth}p{.15\linewidth}p{.15\linewidth}p{.15\linewidth}p{.15\linewidth}}
    \# & \multicolumn{2}{c}{Player} && \multicolumn{2}{c}{Game Master} \\
    \hline

    \theutterance \stepcounter{utterance}  
    & & & \multicolumn{4}{p{0.6\linewidth}}{
        \cellcolor[rgb]{0.9,0.9,0.9}{
            \makecell[{{p{\linewidth}}}]{
                \texttt{\tiny{[P1$\langle$GM]}}
                \texttt{You are an expert in causal inference. The following question is not a typical commonsense query, but rather a meticulously designed question created by a professor specializing in causal inference, intended to assess the students' mastery of the course content.} \\
\texttt{Imagine a self{-}contained, hypothetical world with only the following conditions, and without any unmentioned factors or causal relationships: Wibl has a direct effect on vubr and zuph. Zuph has a direct effect on uvzi. Vubr has a direct effect on uvzi. Method 1: We look at how zuph correlates with uvzi case by case according to vubr. Method 2: We look directly at how zuph correlates with uvzi in general. To understand how zuph affects uvzi, is it more correct to use the Method 1 than Method 2?} \\
\texttt{Start your answer with 'Yes' or 'No', followed by additional reasoning or evidence to support your explanation.} \\
            }
        }
    }
     \\ \\

    \theutterance \stepcounter{utterance}  
    & \multicolumn{4}{p{0.6\linewidth}}{
        \cellcolor[rgb]{0.561,0.737,0.561}{
            \makecell[{{p{\linewidth}}}]{
                \texttt{\tiny{[P1$\rangle$GM]}}
                \texttt{Yes} \\
\\ 
\texttt{To determine the correct approach, we must first map the causal structure described in the hypothetical world. The conditions state that **Wibl** causes both **Vubr** and **Zuph**, making Wibl a common cause (confounder) of the two. Furthermore, **Zuph** has a direct effect on **Uvzi**, and **Vubr** also has a direct effect on **Uvzi**.} \\
\\ 
\texttt{The causal graph is defined by the following edges:} \\
\texttt{1. \$Wibl \textbackslash{}rightarrow Zuph\$} \\
\texttt{2. \$Wibl \textbackslash{}rightarrow Vubr\$} \\
\texttt{3. \$Zuph \textbackslash{}rightarrow Uvzi\$} \\
\texttt{4. \$Vubr \textbackslash{}rightarrow Uvzi\$} \\
\\ 
\texttt{We are interested in estimating the causal effect of **Zuph** on **Uvzi**.} \\
\texttt{The direct causal path is \$Zuph \textbackslash{}rightarrow Uvzi\$.} \\
\texttt{However, there is a backdoor path from Zuph to Uvzi: \$Zuph \textbackslash{}leftarrow Wibl \textbackslash{}rightarrow Vubr \textbackslash{}rightarrow Uvzi\$.} \\
\texttt{This path creates a spurious association between Zuph and Uvzi because Wibl influences both.} \\
\\ 
\texttt{{-} **Method 2** (looking at the general correlation) leaves this backdoor path open. Consequently, the observed correlation will be biased, conflating the true causal effect of Zuph with the confounding influence of Wibl transmitted through Vubr.} \\
\texttt{{-} **Method 1** involves conditioning on **Vubr** (stratifying by Vubr). In the backdoor path \$Zuph \textbackslash{}leftarrow Wibl \textbackslash{}rightarrow Vubr \textbackslash{}rightarrow Uvzi\$, the variable **Vubr** is a non{-}collider (it has an arrow pointing into it from Wibl and an arrow pointing out to Uvzi). Conditioning on a non{-}collider on a backdoor path successfully blocks the flow of spurious association along that path.} \\
\\ 
\texttt{By conditioning on Vubr, we block the backdoor path \$Zuph \textbackslash{}leftarrow Wibl \textbackslash{}rightarrow Vubr \textbackslash{}rightarrow Uvzi\$, thereby isolating the direct causal effect of Zuph on Uvzi. Therefore, Method 1 provides an unbiased estimate of the causal effect, whereas Method 2 does not. Thus, it is more correct to use Method 1.} \\
            }
        }
    }
    & & \\ \\

    \theutterance \stepcounter{utterance}  
    & & & \multicolumn{2}{p{0.3\linewidth}}{
        \cellcolor[rgb]{0.9,0.9,0.9}{
            \makecell[{{p{\linewidth}}}]{
                \texttt{\tiny{[GM$|$GM]}}
                \texttt{Yes} \\
            }
        }
    }
    & & \\ \\

    \theutterance \stepcounter{utterance}  
    & & & \multicolumn{2}{p{0.3\linewidth}}{
        \cellcolor[rgb]{0.9,0.9,0.9}{
            \makecell[{{p{\linewidth}}}]{
                \texttt{\tiny{[GM$|$GM]}}
                \texttt{yes} \\
            }
        }
    }
    & & \\ \\

    \theutterance \stepcounter{utterance}  
    & & & \multicolumn{2}{p{0.3\linewidth}}{
        \cellcolor[rgb]{0.9,0.9,0.9}{
            \makecell[{{p{\linewidth}}}]{
                \texttt{\tiny{[GM$|$GM]}}
                \texttt{game\_result = WIN} \\
            }
        }
    }
    & & \\ \\

\end{supertabular}
}

\end{document}
